\begin{table}[!h]
\centering
\caption{\label{tab:subcohort_did_9a}Sub-Cohort DiD: International Migrant (=1) (Sub-Cohort Decomposition by Developmental Window)}
\centering
\fontsize{9}{11}\selectfont
\begin{threeparttable}
\begin{tabular}[t]{lccc}
\toprule
  & Early childhood
(ages 0–5) & Middle childhood
(ages 6–12) & Adolescence
(ages 13–17)\\
\midrule
\addlinespace[0.6em]
\multicolumn{4}{l}{\textbf{Panel A: High-Conflict Districts defined by Months of War}}\\
\hspace{1em}HC × Sub-Cohort & -5.53*** & -4.66*** & -4.16**\\
\hspace{1em} & (1.51) & (1.56) & (1.63)\\
\hspace{1em}Observations & 21,325 & 19,403 & \vphantom{1} 15,906\\
\addlinespace[0.6em]
\multicolumn{4}{l}{\textbf{Panel B: High-Conflict Districts defined by Casualties}}\\
\hspace{1em}HC × Sub-Cohort & -6.38*** & -6.01*** & -4.65***\\
\hspace{1em} & (1.40) & (1.40) & (1.66)\\
\hspace{1em}Observations & 21,325 & 19,403 & 15,906\\
\bottomrule
\end{tabular}
\begin{tablenotes}
\item \textit{Notes:} 
\item Outcome: International Migrant (=1), scaled by 100.
\item Each cell reports the DiD coefficient (HC × Sub-Cohort) from a separate regression.
\item Control group is held constant: pooled individuals aged 26–40 at conflict onset (C: 26-35 + C: 36-40; currently 47–61 in 2017).
\item All specifications include district fixed effects, age fixed effects, and demographic controls (sex).
\item Standard errors in parentheses, clustered at district level.
\item *** p<0.01, ** p<0.05, * p<0.10.
\item Source: Nepal Labour Force Survey 2017/18; INSEC.
\end{tablenotes}
\end{threeparttable}
\end{table}
